\documentclass[a4paper, 12p, titlepaget]{book}
\usepackage{cmbright}
\usepackage{graphicx}
\usepackage[intlimits]{amsmath}
\usepackage{amssymb}
\usepackage{amsthm}
\usepackage{parskip}
\usepackage{pgf}
\usepackage{tikz}
\usetikzlibrary{shapes}
\usetikzlibrary{positioning}
\usepackage{enumerate}
\usepackage{framed}

\newcommand \name {NESIZER }
\newcommand {\btn}[1] {\framebox{\textbf{\footnotesize #1}}}
\newcommand {\cbtn}[1] {\tikz \node[draw,circle]{\textbf{\footnotesize #1}};}
\newcommand {\lbl}[1] {\emph{\footnotesize #1}}
\newcommand \dsp \texttt


\title{NESIZER}
\author{}
\date{}


\begin{document}

\definecolor{buttonredh}{RGB}{255,0,0}
\definecolor{buttonred}{RGB}{240, 161, 161}
\definecolor{buttongray}{RGB}{71,71,71}

\tikzstyle{mbutton}=[font=\tiny,rectangle, draw, fill=buttonred, text=gray, minimum height=1.2cm, minimum width=1.2cm]
\tikzstyle{mhbutton}=[font=\tiny,rectangle, draw, fill=buttonredh, minimum height=1.2cm,minimum width=1.2cm]

\tikzstyle{hbutton}=[font=\tiny,rectangle, draw, fill=black, text=white, minimum height=1.2cm, minimum width=1.2cm]
\tikzstyle{button}=[font=\tiny,rectangle, draw, fill=buttongray, text=gray, minimum height=1.2cm,minimum width=1.2cm]

\tikzstyle{dummybutton}=[font=\tiny,rectangle, fill=white, text=white, minimum height=1cm,minimum width=1cm]
\tikzstyle{hcircbutton}=[font=\tiny,circle, fill=black, text=white, minimum height=1.2cm,minimum width=1.2cm]
\tikzstyle{circbutton}=[font=\tiny,circle, fill=buttongray, text=gray, minimum height=1.2cm,minimum width=1.2cm]


\begin{titlepage}
  \begin{center}
    \vfil
    \includegraphics[width=12cm]{../nesizer_black.png}\\[3cm]

    \textsc{\LARGE{8-bit synthesizer}}\\[8cm]

    \textsc{\Large{Operating manual}}
  \end{center}
\end{titlepage}

\tableofcontents

%\maketitle

\chapter{User manual}

\section{Overview}

\section{The NESIZER}

\subsection{Sound channels}

At the heart of the \name is the NES APU chip, usually called 2A03 (or 2A07 if you use a chip from the PAL NES). The APU has five separate sound channels:

\begin{itemize}
\item \textbf{SQ1} and \textbf{SQ2}: These produce square waves with three selectable \emph{duty cycles} (pulse widths).
\item \textbf{TRI}: This channel produces triangular waves, but with a low 4 bit amplitude resolution. This results in the characteristic aliased NES bass and flute tones.
\item \textbf{NOISE}: This channel produces various forms of noise. The noise can be white noise, or pitched noise if the \lbl{LOOP} mode is engaged.
\item \textbf{DMC}: This channel can output 7-bit samples at a quick rate. The sampling rate is 16 kHz in the \name.
\end{itemize}

\subsection{Modulation}

Because the 2A03 is put under much tighter control in the \name than in a NES or Famicom console, the \name features extensive modulation capabilities.

\begin{itemize}
\item Three separate low frequency oscillators \textbf{LFO1}, \textbf{LFO2}, \textbf{LFO3} with selectable waveforms (ramp up, ramp down, sine wave, triangle wave or square wave)
\item Dedicated ADSR envelope generators for the square and noise channels
\item Portamento / glide for the square and triangle channels
\end{itemize}

\section{PROGRAM and PAGE 2 modes}

This is the active mode when the \name starts up. When in other modes, press the \btn{PROGRAM} button to switch to this mode.

In the programming mode, the various parameters of the sound channels and LFOs can be changed. A set of channel and LFO settings are collectively known as a \emph{patch}. The \name can save patches to memory, and has room for 100 patches.

\begin{tikzpicture}[align=center, node distance=1.5cm, auto]
  \node [mhbutton] (1) {SQ1};
  \node [mhbutton, right of=1] (2) {SQ2};
  \node [mhbutton, right of=2] (3) {TRI};
  \node [mhbutton, right of=3] (4) {NOISE};
  \node [mhbutton, right of=4] (5) {DMC};
  \node [hbutton, right of=5] (6) {LFO1};
  \node [hbutton, right of=6] (7) {LFO2};
  \node [hbutton, right of=7] (8) {LFO3};

  \node [hbutton, below=1] (9) {};
  \node [hbutton, right of=9] (10) {};
  \node [hbutton, right of=10] (11) {};
  \node [hbutton, right of=11] (12) {};
  \node [hbutton, right of=12] (13) {};
  \node [hbutton, right of=13] (14) {};
  \node [hbutton, right of=14] (15) {};
  \node [hbutton, right of=15] (16) {};
\end{tikzpicture}

\begin{tikzpicture}[align=center, node distance=1.5cm, auto]
  \node [dummybutton] (1) {};
  \node [dummybutton, right of=1] (2) {};
  \node [hcircbutton, right of=2] (3) {PROGRAM};
  \node [circbutton, right of=3] (4) {PAGE 2};
  \node [circbutton, right of=4] (5) {ASSIGNER};
  \node [circbutton, right of=5] (6) {SETTINGS};
  \node [dummybutton, right of=6] (7) {};
  \node [dummybutton, right of=7] (8) {};

  \node [hcircbutton, below=1] (8) {CLEAR};
  \node [hcircbutton, right of=8] (9) {SAVE};
  \node [dummybutton, right of=9] (10) {};
  \node [dummybutton, right of=10] (11) {};
  \node [dummybutton, right of=11] (12) {};
  \node [dummybutton, right of=12] (13) {};
  \node [hcircbutton, right of=13] (14) {UP};
  \node [hcircbutton, right of=14] (15) {DOWN};
\end{tikzpicture}

\subsection{Loading patches}
To load a patch, press either \btn{UP} or \btn{DOWN}. To go quickly up or down, press and hold the respective button. The current patch number is shown on the numeric display.

\subsection{Saving patches}
To save a patch, press \btn{SAVE}. The button will start to blink to indicate that you can select where to store the new patch. Use \btn{UP} and \btn{DOWN} to select where to store the patch. Press \btn{SAVE} again to store the patch at the selected location.

\emph{Note: When channel and LFO settings are changed, these are not saved until you press \btn{SAVE}.}

\subsection{Enabling and disabling channels}
To enable or disable a channel, press the corresponding channel button. When a channel is disabled, it does not produce sound when being triggered by the sequencer or incoming MIDI data.

\subsection{PAGE 2}

By pressing \btn{PAGE 2}, a second ``page'' og parameters become available. Some of them are channel parameters, and others are related to note assignment. When Page 2 is selected, the \name is still in programmer mode. The parameters that are avilable in page 2 are shown in yellow color in the figures below.

\subsection{Changing channel parameters}
To change a channel's parameter, press and hold the desired channel's button, and the desired parameter button. For example, to change the attack of the square 1 channel, press \btn{SQ1} and \btn{A}. The button LEDs will start to blink to indicate which channel parameter is being changed. Use the \btn{UP} and \btn{DOWN} buttons to change the parameter value. When you have the desired value, press \btn{SAVE}.

\subsubsection{Square channels}

\begin{tikzpicture}[align=center, node distance=1.5cm and 0.5cm]
  \node [mhbutton] (1) {SQ1};
  \node [mhbutton, right of=1] (2) {SQ2};
  \node [mbutton, right of=2] (3) {TRI};
  \node [mbutton, right of=3] (4) {NOISE};
  \node [mbutton, right of=4] (5) {DMC};
  \node [hbutton, right of=5] (6) {LFO1\\[4]\textcolor{yellow}{BEND}};
  \node [hbutton, right of=6] (7) {LFO2\\[4]\textcolor{yellow}{ENV MOD}};
  \node [hbutton, right of=7] (8) {LFO3\\[4]\textcolor{yellow}{VOL MOD}};

  \node [hbutton, below=1] (9) {DUTY};
  \node [hbutton, right of=9] (10) {DETUNE};
  \node [hbutton, right of=10] (11) {OCTAVE};
  \node [hbutton, right of=11] (12) {GLIDE};
  \node [hbutton, right of=12] (13) {A};
  \node [hbutton, right of=13] (14) {D};
  \node [hbutton, right of=14] (15) {S};
  \node [hbutton, right of=15] (16) {R};
\end{tikzpicture}

\begin{tabular}{l | l | l}
  Parameter & Description & Range\\ \hline
  \textbf{LFO1} & Intensity of modulation by LFO1 & 0 - 99\\
  \textbf{LFO2} & Intensity of modulation by LFO2 & 0 - 99\\
  \textbf{LFO3} & Intensity of modulation by LFO3 & 0 - 99\\
  \textbf{DUTY} & Duty cycle & 0 - 3\\
  \textbf{DETUNE} & Detuning & -9 - 9\\
  \textbf{OCTAVE} & Octave & -4 - 4\\
  \textbf{GLIDE} & Portamento glide time & 0 - 99\\
  \textbf{A} & Volume envelope attack & 0 - 99\\
  \textbf{D} & Volume envelope decay & 0 - 99\\
  \textbf{S} & Volume envelope sustain & 0 - 15\\
  \textbf{R} & Volume envelope release & 0 - 99\\
  \textbf{BEND} (page 2) & Bend wheel intensity in semitones & 0 - 24\\
  \textbf{ENV MOD} (page 2) & Pitch envelope modulation amount & -9 - 9\\
  \textbf{VOL MOD} (page 2) & Volume modulation by LFO3 & 0 - 16\\
\end{tabular}

\subsubsection{Triangle channel}

\begin{tikzpicture}[align=center, node distance=1.5cm and 0.5cm]
  \node [mbutton] (1) {SQ1};
  \node [mbutton, right of=1] (2) {SQ2};
  \node [mhbutton, right of=2] (3) {TRI};
  \node [mbutton, right of=3] (4) {NOISE};
  \node [mbutton, right of=4] (5) {DMC};
  \node [hbutton, right of=5] (6) {LFO1\\[4]\textcolor{yellow}{BEND}};
  \node [hbutton, right of=6] (7) {LFO2\\[4]\textcolor{yellow}{ENV MOD}};
  \node [hbutton, right of=7] (8) {LFO3};

  \node [button, below=1] (9) {};
  \node [hbutton, right of=9] (10) {DETUNE};
  \node [hbutton, right of=10] (11) {OCTAVE};
  \node [hbutton, right of=11] (12) {GLIDE};
  \node [button, right of=12] (13) {};
  \node [button, right of=13] (14) {};
  \node [button, right of=14] (15) {};
  \node [button, right of=15] (16) {};
\end{tikzpicture}

\begin{tabular}{l | l | l}
  Parameter & Description & Range\\ \hline
  \textbf{LFO1} & Intensity of modulation by LFO1 & 0 - 99\\
  \textbf{LFO2} & Intensity of modulation by LFO2 & 0 - 99\\
  \textbf{LFO3} & Intensity of modulation by LFO3 & 0 - 99\\
  \textbf{DETUNE} & Detuning & -9 - 9\\
  \textbf{OCTAVE} & Octave & -4 - 4\\
  \textbf{GLIDE} & Portamento glide time & 0 - 99\\
  \textbf{BEND} (page 2) & Bend wheel intensity in semitones & 0 - 24\\
  \textbf{ENV MOD} (page 2) & Pitch envelope modulation (noise envelope) & -9 - 9\\
\end{tabular}

\subsubsection{Noise channel}

\begin{tikzpicture}[align=center, node distance=1.5cm and 0.5cm]
  \node [mbutton] (1) {SQ1};
  \node [mbutton, right of=1] (2) {SQ2};
  \node [mbutton, right of=2] (3) {TRI};
  \node [mhbutton, right of=3] (4) {NOISE};
  \node [mbutton, right of=4] (5) {DMC};
  \node [hbutton, right of=5] (6) {LFO1\\[4]\textcolor{yellow}{BEND}};
  \node [hbutton, right of=6] (7) {LFO2\\[4]\textcolor{yellow}{ENV MOD}};
  \node [hbutton, right of=7] (8) {LFO3\\[4]\textcolor{yellow}{VOL MOD}};

  \node [hbutton, below=1] (9) {LOOP};
  \node [button, right of=9] (10) {};
  \node [button, right of=10] (11) {};
  \node [button, right of=11] (12) {};
  \node [hbutton, right of=12] (13) {A};
  \node [hbutton, right of=13] (14) {D};
  \node [hbutton, right of=14] (15) {S};
  \node [hbutton, right of=15] (16) {R};
\end{tikzpicture}

\begin{tabular}{l | l | l}
  Parameter & Description & Range\\ \hline
  \textbf{LFO1} & Intensity of modulation by LFO1 & 0 - 99\\
  \textbf{LFO2} & Intensity of modulation by LFO2 & 0 - 99\\
  \textbf{LFO3} & Intensity of modulation by LFO3 & 0 - 99\\
  \textbf{LOOP} & Looped noise & on/off\\
  \textbf{A} & Volume envelope attack & 0 - 99\\
  \textbf{D} & Volume envelope decay & 0 - 99\\
  \textbf{S} & Volume envelope sustain & 0 - 15\\
  \textbf{R} & Volume envelope release & 0 - 99\\
  \textbf{BEND} (page 2) & Bend wheel intensity (in steps) & 0 - 15\\
  \textbf{ENV MOD} (page 2) & Pitch envelope modulation amount & -9 - 9\\
  \textbf{VOL MOD} (page 2) & Volume modulation by LFO3 & 0 - 16\\
\end{tabular}

\subsubsection{DMC channel}

\begin{tikzpicture}[align=center, node distance=1.5cm and 0.5cm]
  \node [mbutton] (1) {SQ1};
  \node [mbutton, right of=1] (2) {SQ2};
  \node [mbutton, right of=2] (3) {TRI};
  \node [mbutton, right of=3] (4) {NOISE};
  \node [mhbutton, right of=4] (5) {DMC};
  \node [button, right of=5] (6) {LFO1};
  \node [button, right of=6] (7) {LFO2};
  \node [button, right of=7] (8) {LFO3};

  \node [hbutton, below=1] (9) {LOOP};
  \node [hbutton, right of=9] (10) {DIVIDER};
  \node [button, right of=10] (11) {};
  \node [button, right of=11] (12) {};
  \node [button, right of=12] (13) {};
  \node [button, right of=13] (14) {};
  \node [button, right of=14] (15) {};
  \node [button, right of=15] (16) {};
\end{tikzpicture}

\begin{tabular}{l | l | l}
  Parameter & Description & Range\\ \hline
  \textbf{LOOP} & Play samples in loop & on/off\\
  %\textbf{DIVIDER} & Divider for sample rate & 1 - 4\\
\end{tabular}

\subsubsection{LFOs}

\begin{tikzpicture}[align=center, node distance=1.5cm and 0.5cm]
  \node [mbutton] (1) {SQ1};
  \node [mbutton, right of=1] (2) {SQ2};
  \node [mbutton, right of=2] (3) {TRI};
  \node [mbutton, right of=3] (4) {NOISE};
  \node [mbutton, right of=4] (5) {DMC};
  \node [hbutton, right of=5] (6) {LFO1};
  \node [hbutton, right of=6] (7) {LFO2};
  \node [hbutton, right of=7] (8) {LFO3};

  \node [hbutton, below=1] (9) {WAVE};
  \node [button, right of=9] (10) {};
  \node [button, right of=10] (11) {};
  \node [button, right of=11] (12) {};
  \node [button, right of=12] (13) {};
  \node [button, right of=13] (14) {};
  \node [button, right of=14] (15) {};
  \node [button, right of=15] (16) {};
\end{tikzpicture}

\begin{tabular}{l | l | l}
  Parameter & Description & Range\\ \hline
  \textbf{WAVE} & LFO waveform & 1 - sine\\
  & & 2 - ramp up\\
  & & 3 - ramp down\\
  & & 4 - square\\
  & & 5 - triangle
\end{tabular}


\section{SETTINGS}

When in the settings mode, various aspects of how the \name operates can be changed. In this mode, the buttons have the following functions:

\begin{tikzpicture}[align=center, node distance=1.5cm]
  \node [mhbutton] (1) {SQ1};
  \node [mhbutton, right of=1] (2) {SQ2};
  \node [mhbutton, right of=2] (3) {TRI};
  \node [mhbutton, right of=3] (4) {NOISE};
  \node [mhbutton, right of=4] (5) {DMC};
  \node [hbutton, right of=5] (6) {MIDI CHN};
  \node [hbutton, right of=6] (7) {};
  \node [hbutton, right of=7] (8) {CHIP TYPE};

  \node [hbutton, below=1] (9) {PATCH\\RESET};
  \node [button, right of=9] (10) {};
  \node [button, right of=10] (11) {};
  \node [button, right of=11] (12) {};
  \node [button, right of=12] (13) {};
  \node [button, right of=13] (14) {};
  \node [hbutton, right of=14] (15) {SAMPLE\\RESET};
  \node [hbutton, right of=15] (16) {SAMPLE\\DELETE};

  \node [circbutton, below= of 9] (17) {};
  \node [circbutton, right of=17] (18) {};
  \node [dummybutton, right of=18] (19) {};
  \node [dummybutton, right of=19] (20) {};
  \node [dummybutton, right of=20] (21) {};
  \node [dummybutton, right of=21] (22) {};
  \node [hcircbutton, right of=22] (23) {UP};
  \node [hcircbutton, right of=23] (24) {DOWN};
\end{tikzpicture}

\btn{UP} and \btn{DOWN} can be used to select a sample number. Sample numbers that are occupied are marked with a dot on the display. A selected sample can be deleted by pressing \btn{DELETE}.

\subsection{MIDI}

The \name can be controlled externally using MIDI.

\subsubsection{Assigning MIDI channels}

Each of the five sound channels can be assigned to any of the 16 MIDI channels, and will then only respond to incoming messages on the selected channel.

To select a MIDI channel, enter the \lbl{SETTINGS} mode by pressing \btn{SETTINGS}. Hold down \btn{MIDI CHANNEL} and then press the desired channel's button. The LEDs will flash and you can select one of the 16 MIDI channels using the \btn{UP} and \btn{DOWN} buttons. If you select the value 0, the sound channel will not listen to any MIDI channel.

\subsection{Checking the 2A03 type}
To see which 2A03 chip is being used, press the LFO3 button. One of the following numbers will show up:
\begin{itemize}
\item 12: Genuine RP2A03
\item 15: 2A03 clone, e.g. 6527P
\item 16: 2A07 (PAL version of 2A03)
\end{itemize}

\subsection{Resetting patches}

Press \btn{PATCH RESET} to delete all patches. Every patch will be initialized to a basic patch with no channels enabled, square duty cycles set at 50\% and full envelope sustain levels with no attack, decay or release.

\subsection{Maintaining samples}

When in SETTINGS mode, the up and down buttons are used to select DMC samples. A dot appearing on the display indicates that the selected sample location is occupied. When an occupied sample is selected, it can be deleted by pressing \btn{SAMPLE DELETE}. All samples can be erased by pressing \btn{SAMPLE RESET}.


\end{document}
